\chapter{Qadar and the Resolution of Determinism}
\label{ch:qadar}

\epigraph{\itshape Indeed, all things We created with predestination
(\arb{qadar}).}{--- \Quran{54:49}}

\noindent
This chapter treats the most theologically deep of the worked
treatments. The Islamic doctrine of \arb{qadar} (divine decree)
appears, on the standard reading, to be in tension with the
parallel doctrine of \arb{ikhtiyar} (human choice) and with the
moral responsibility that the same Quran insists upon. The classical
\arb{kalam} debate (Mu`tazila, Ash`ariyya, Maturidiyya) attempted
to articulate positions that preserve both. Each position has its
strengths; none has produced a fully satisfying resolution.

The chapter argues that the apparent tension is the artifact of
Newtonian assumptions about determinism that quantum mechanics has
shown to be inappropriate. In the Bohmian and QBist frameworks of
\xref{ch:quantum-foundations}, determinism and the reality of
motion can both hold; probability can be irreducibly perspectival
without contradiction. Either framework provides the structural
form in which qadar and ikhtiyar are coherent.

The chapter is Level 2 in our taxonomy. The classical theological
problem is structurally dissolved; the metaphysical content is
elucidated; the operational mechanism (in the strong scientific
sense) is not specified.

The chapter proceeds in eight sections. \xref{sec:qad-claim}
states the doctrine. \xref{sec:qad-classical} surveys the classical
positions. \xref{sec:qad-newtonian} identifies the Newtonian source
of the apparent contradiction. \xref{sec:qad-bohmian-resolution}
provides the Bohmian resolution. \xref{sec:qad-qbism-resolution}
provides the QBist resolution. \xref{sec:qad-akbarian} integrates
the Akbarian framework. \xref{sec:qad-economic} draws the economic
implications. \xref{sec:qad-implications} draws the broader
implications.

\section{The doctrine}
\label{sec:qad-claim}

\subsection{The Quranic loci}

The Quranic doctrine of \arb{qadar} is comprehensive. \Quran{54:49}
asserts: ``All things We created with predestination.'' \Quran{57:22}:
``No disaster strikes upon the earth or among yourselves except
that it is in a register before We bring it into being.'' \Quran{76:30}
and \Quran{81:29}: ``You do not will except that Allah wills.''

The loci establish that the divine decree extends over all events,
including the events that human action produces. The decree is not
limited to natural events; it includes the willed actions of
created agents. This is the strong form of the doctrine.

\subsection{The complementary doctrine of choice}

The same Quran insists that humans are responsible for their
actions and will be judged accordingly. \Quran{18:29}: ``Whoever
wills, let him believe; whoever wills, let him disbelieve.''
\Quran{2:286}: ``Allah does not burden a soul beyond that it can
bear.'' \Quran{91:7-10}: ``By the soul and He who proportioned it
and inspired it with discernment of its wickedness and its
righteousness, he has succeeded who purifies it.''

Moral responsibility presupposes meaningful choice. The believer
is responsible for choosing belief and for the consequent
behaviors; the disbeliever is responsible for choosing disbelief.
Neither responsibility makes sense if the choice is illusory.

\subsection{The apparent contradiction}

The two doctrines together produce the apparent contradiction. If
all events are determined by divine decree, then human choice is
either illusory (the decree determines what will be willed) or
the decree does not extend to human action (in which case
\Quran{54:49} is not in its strongest form). Either horn of the
dilemma sacrifices something the tradition holds essential.

The classical \arb{kalam} debate is the history of attempts to
articulate a coherent position. We treated this in
\xref{sec:kalam}; we recall the major positions briefly here.

\section{The classical positions}
\label{sec:qad-classical}

\subsection{Mu`tazili libertarianism}

The Mu`tazila held that humans create their own actions. The
divine decree extends to natural events but not to the willed
actions of agents. This preserves moral responsibility (humans
are genuine causes of their actions) but at the cost of
restricting divine creative activity.

The position has the strength of clarity. Moral responsibility
makes sense; choice is real; agents are causes. The cost is the
restriction of divine creative activity, which the Quranic doctrine
appears to assert without restriction.

\subsection{Ash`arite kasb}

The Ash`ariyya held that all events, including human actions, are
created by Allah. The doctrine of \arb{kasb} (acquisition) was
developed to preserve human responsibility: humans ``acquire''
their actions through some mode of involvement, even though the
actions are created by Allah.

The position preserves divine omnipotence (everything is divinely
created) but the doctrine of kasb is famously difficult to
articulate clearly. What exactly is acquisition, if the action
itself is divinely created? The literature on the topic is
extensive but contested.

\subsection{Maturidi compromise}

The Maturidiyya distinguished between the capacity for action
(divinely created) and the exercise of the capacity (the human's
own). The compromise has the merit of distinguishing different
ontological levels at which divine and human causality operate;
it has the cost of inheriting criticism from both sides.

\subsection{The unresolved tension}

None of the classical positions produces a fully satisfying
resolution. Each addresses the problem within the conceptual
apparatus available to its formulators; each leaves residual
puzzles.

The puzzle is not just an abstract theological one. It bears on
practical questions: how should the believer engage with their own
choices? What does it mean to strive when outcomes are decreed?
How does the believer reconcile responsibility with the
recognition that Allah controls all events?

\section{The Newtonian source of the apparent contradiction}
\label{sec:qad-newtonian}

\subsection{Newtonian determinism}

The classical theological discussion presupposed a particular
conception of determinism. In the Newtonian conception, a
deterministic system is one whose future is fully fixed by its
present state and the laws of motion. Given the initial conditions
and the dynamics, the future trajectory is uniquely determined.

In this conception, motion is the unfolding of what was already
implicit in the initial conditions. The present state contains the
future trajectory; the trajectory is just the explication of the
present. There is no genuine novelty; nothing new happens; the
future is just the temporal unrolling of the present.

\subsection{Why this conception generates the contradiction}

Under this conception, if the divine decree fixes the initial
conditions (as the doctrine of qadar implies), then the future
trajectory is fixed. Human ``choice'' is just the temporal
unrolling of what the initial conditions already determined. The
choice is not genuine; the agent is not really doing anything new;
the moral responsibility loses its grip.

This is the structural source of the apparent contradiction. The
Newtonian conception of determinism makes determination and
genuine motion mutually exclusive. The classical theologians,
working with this conception (or something like it), could not
both preserve determinism and preserve genuine motion.

\subsection{The Newtonian conception is wrong}

The Newtonian conception of determinism has been superseded by the
twentieth-century discoveries we treated in
\xref{ch:quantum-foundations}. Quantum mechanics has shown that
the world is not Newtonian; the relationship between
determination and motion is not what classical mechanics
assumed.

In particular, two interpretations of quantum mechanics ---
Bohmian mechanics and QBism --- exhibit structures in which
determinism and the reality of motion can both hold. We use these
structures to dissolve the apparent contradiction.

\section{The Bohmian resolution}
\label{sec:qad-bohmian-resolution}

\subsection{The Bohmian structure}

We developed Bohmian mechanics in \xref{sec:interpretations} and
its application to the rizq-conservation claim in
\xref{sec:rizq-bohmian}. We recall the key features.

In Bohmian mechanics, a particle has a definite position at all
times. The position evolves according to a guidance equation
involving the wave function. The wave function evolves
deterministically (Schr\"odinger equation). The trajectory of the
particle is fully determined by the initial wave function and the
particle's initial position. Yet the motion is real; the particle
actually moves.

This is the structural form we need. The trajectory is determined;
the motion is real. There is no contradiction.

\subsection{Application to qadar}

\begin{conceptbox}[Bohmian formulation of qadar]
\label{conc:bohmian-qadar}
The configuration space for human actions is the space of all
possible action sequences. The wave function $\Psi$ on this
space is determined by the divine decree. The actual life of the
agent is a single trajectory through configuration space, guided
by $\Psi$.

The trajectory is fully determined by $\Psi$ and the agent's
initial position. The motion along the trajectory is real; the
agent actually performs the actions. Both descriptions are
correct; both are required.
\end{conceptbox}

The Bohmian formulation provides exactly the structural resolution
the classical theologians were seeking. The Ash`arite intuition
that everything is divinely created is preserved (the wave
function is the divine decree). The Mu`tazili intuition that
human action is real is preserved (the motion along the trajectory
is real). The two are not in tension; they are descriptions of
different aspects of the same structural reality.

\subsection{What this resolution accomplishes}

The Bohmian resolution accomplishes several things.

First, it dissolves the apparent contradiction. The classical
worry that determinism makes motion illusory was a worry about a
structure that, in Bohmian mechanics, does not obtain. The
trajectory can be fully determined and the motion can be fully
real; the structure permits both.

Second, it explains why the classical positions had merit. The
Mu`tazili were right that human action is real; the Ash`ariyya
were right that the wave function (the structural ground) is
divinely determined. Each position captured part of the truth.

Third, it provides a framework for thinking about practical
questions. The agent's striving is the motion along the
trajectory; it is real and effective in producing the trajectory.
The trajectory itself is determined; the agent is responsible
for being the kind of agent whose trajectory has the character it
has.

\subsection{The structural nature of the resolution}

The Bohmian resolution is structural, not literal. We are not
claiming that the human mind is governed by Bohmian quantum
dynamics. We are claiming that the structure ``determined
trajectory + real motion'' is mathematically coherent and
physically instantiated in Bohmian mechanics. Whatever the actual
mechanism of human action is, it can have this structure without
contradiction.

This is Level 2 in our taxonomy. The structural problem is
dissolved; the operational mechanism in the strong scientific
sense is not specified. The dissolution is what was needed; the
specification of the underlying mechanism is a different and
harder question.

\section{The QBist resolution}
\label{sec:qad-qbism-resolution}

\subsection{The QBist structure}

QBism (Quantum Bayesianism) was treated in \xref{sec:interpretations}.
We recall the key feature: probability in QBism is irreducibly
perspectival. Different agents, with different information, can
have different valid probability assignments for the same event.
There is no single objective probability; probability is the
relationship between an agent and an event.

\subsection{Application to qadar and ikhtiyar}

\begin{conceptbox}[QBist formulation of qadar]
\label{conc:qbist-qadar}
The probability of a future event is perspectival. From the divine
perspective, knowledge of all events is perfect; the divine
``probability'' is concentrated on the actual outcome
(probability 1 for what will happen, probability 0 for what will
not). From the human perspective, knowledge is finite; the human
``probability'' is spread over multiple possible outcomes.

Both probability assignments are correct. They reflect the
different epistemic positions of the divine and human
perspectives. There is no contradiction.
\end{conceptbox}

The QBist resolution complements the Bohmian one. Where Bohmian
mechanics gives us the structure of determinism plus real motion,
QBism gives us the structure of perspectival probability. Either
suffices for dissolving the classical contradiction; together they
reinforce each other.

\subsection{The economic implications of QBism}

QBism has direct implications for economic decision-making. The
agent's decisions should be based on the agent's probability
assignment, which is the agent's epistemic state, not on some
objective probability that the agent cannot access. This is
already the standard practice in modern decision theory; QBism
provides the metaphysical justification.

The believer's probability assignments are informed by their
reliance on Allah's knowledge of the future, but the believer
themselves does not have direct access to that knowledge. The
believer therefore makes decisions based on their finite
epistemic state, with appropriate trust in the divine knowledge
that exceeds their own. This is the structure of \arb{tawakkul}
(trust): the believer relies on Allah while still acting on the
basis of their own finite knowledge.

\section{The Akbarian integration}
\label{sec:qad-akbarian}

\subsection{The unity of being and the multiplicity of perspectives}

The Akbarian framework of \xref{ch:akbarian} provides the
ontological substrate. The doctrine of \arb{wahdat al-wujud}
asserts that there is one ultimate reality (\arb{wujud}); all
created things are loci of divine self-disclosure. The unity of
\arb{wujud} corresponds to the perfect divine knowledge of all
events. The multiplicity of loci corresponds to the finite
perspectives of the dissociated alters.

In this framework, the divine perspective and the human
perspective are not two separate things in tension. They are the
same underlying reality (\arb{wujud}) viewed from two positions
(the unified source vs. the dissociated alter). Both perspectives
are real; both are correct; the apparent tension is the artifact
of treating them as competing descriptions of the same reality
rather than as complementary aspects of a single reality.

\subsection{Tajalli and the dynamics of decree}

The Akbarian doctrine of \arb{tajalli} (continuous self-disclosure,
\xref{sec:tajalli-mechanism}) gives the dynamic content. The
divine decree is not a static arrangement; it is the continuous
self-disclosure of \arb{wujud} in the loci. Each moment is a
fresh manifestation; the trajectory of each soul is the temporal
sequence of these manifestations.

The agent's striving is the agent's participation in the
manifestation. The agent does not create the manifestation
(divine self-disclosure is the source); but the agent's own
structural state determines what manifestation is appropriate.
The agent who is properly oriented receives a manifestation that
unfolds appropriately; the agent who is misoriented receives a
manifestation that unfolds in distorted form.

This is the deep content of the practical injunction to
\arb{tawakkul}. The believer is asked to orient themselves
properly toward the divine self-disclosure. The orientation does
not change the underlying decree (which is determined by the
ontological structure); the orientation determines how the decree
is experienced.

\subsection{Insan al-kamil and the perfect alignment}

The Akbarian doctrine of \arb{insan al-kamil} (the perfect human,
\xref{sec:perfect-human}) is the limiting case. The perfect human
is the locus that manifests all the divine names in equilibrium;
their trajectory through configuration space is the perfectly
aligned manifestation of the underlying decree. The Prophet, on
the Akbarian reading, is the exemplar.

For ordinary believers, the goal is approximation: to align
oneself as fully as possible with the divine names so that one's
trajectory is as close as possible to the perfect manifestation.
The practical Islamic life --- prayer, fasting, charity, the
prohibitions and commandments --- is the structural protocol for
this alignment.

\section{The economic implications}
\label{sec:qad-economic}

\subsection{Why qadar matters for economic theory}

The doctrine of qadar has direct economic implications. The
believer who properly understands the doctrine engages with
economic activity differently from the agent who does not.

\textbf{Effort under qadar.} The agent's striving is real and
effective; it is the agent's participation in the manifestation
of the trajectory. But the trajectory's outcome is not the
agent's to determine; the outcome belongs to the underlying
decree. The agent is therefore committed to effort but not to
outcome. The classical teaching of ``tie your camel and trust in
Allah'' (\hadith{Tirmidhi} 2517) captures this practical
orientation.

\textbf{Risk under qadar.} The agent does not bear risk in the
ultimate sense; the outcome is determined by the decree, not by
the agent's risk-taking. But the agent's risk-taking is real
within the trajectory; the manifestation includes the agent's
risk-bearing. This permits the equity-participation contracts of
Islamic finance to be coherent: the agent shares the risk as
participation in the manifestation, not as bearing of cosmic
responsibility for outcome.

\textbf{Provision under qadar.} The agent's provision is determined
by the decree (per \xref{ch:rizq-conservation}); the agent's
striving determines how it manifests temporally. The proper
attitude toward provision is therefore neither anxious effort to
maximize nor passive resignation; it is engaged participation in
the manifestation of the decreed provision.

\subsection{The connection to ergodicity}

The qadar framework illuminates the ergodicity discussion of
\xref{ch:ergodicity}. The ensemble average of wealth corresponds,
in qadar terms, to the average across counterfactual trajectories.
The time average corresponds to the actual trajectory. The
agent experiences the time average, not the ensemble average;
this is what corresponds to the agent's actual experience of
their decreed life.

The ergodicity-economics result that the time average differs
from the ensemble average has a parallel in qadar terms. The
agent who optimizes for the ensemble average (acting as if the
average across counterfactuals were the relevant criterion) makes
worse decisions than the agent who optimizes for the time average
(acting as if the actual trajectory were the relevant criterion).
The traditional Islamic guidance to \arb{tawakkul} is, in part,
the practical instruction to optimize for the time average rather
than the ensemble average.

\section{Broader implications}
\label{sec:qad-implications}

\subsection{For the classical kalam debate}

The chapter has implications for how we read the classical
\arb{kalam} debate. The debate was substantively serious; the
classical theologians were addressing a real problem. The
problem, however, was unresolvable within the conceptual
apparatus available to them. The Newtonian framework that was
implicit in their thinking made determinism and motion mutually
exclusive in the way that quantum mechanics has shown to be
unnecessary.

This is not a criticism of the classical theologians. They were
working with the conceptual resources available in their time and
they did remarkable work within those constraints. The modern
resources we have available --- Bohmian mechanics, QBism --- were
not available to them. Our resolution of the contradiction is
possible because we have these resources; the resolution does
not refute the classical positions but extends them with new
conceptual tools.

\subsection{For modern philosophical theology}

The chapter has implications for the broader contemporary
discussion of free will and determinism. The libertarian and
compatibilist positions in modern philosophy have been working
within frameworks similar to the classical kalam ones; they have
generated similar puzzles for similar reasons. The structural
resolution offered here applies broadly.

The implication is that the contemporary free-will/determinism
debate may be addressing a problem whose conceptual basis has
been superseded. The proper question is not ``is determinism
compatible with free will'' (the question that assumes the
Newtonian conception); the proper question is ``what structural
form of determination is consistent with what observable phenomena
of choice.'' The Bohmian and QBist structures provide candidate
answers.

\subsection{For practical religious life}

The chapter has implications for practical religious life. The
believer's relationship with the divine decree is structurally
clarified. The decree is not a pre-existing arrangement that
makes the believer's choices meaningless; it is the
multidimensional reality within which the believer's choices are
the manifestation of the decree. The believer is responsible for
being the agent whose choices manifest the decree well; the
outcome of the manifestation is not the believer's to determine.

This is a clarification, not a replacement, of the traditional
practical guidance. The traditional guidance to strive while
trusting, to plan while accepting, to take risks while
recognizing the limits of one's knowledge --- this guidance is
correct and is now structurally explained.

\section{Conclusion}

Chapter 16 has dissolved the apparent contradiction between qadar
and ikhtiyar by showing that the contradiction was the artifact of
Newtonian assumptions about determinism that quantum mechanics has
made obsolete. The Bohmian and QBist structures provide
mathematical frameworks in which determination and the reality of
motion can both hold; perspectival probability removes the
apparent tension between divine perfect knowledge and human
finite knowledge.

The resolution is Level 2: the structural problem is dissolved;
the operational mechanism in the strong scientific sense is not
specified. The classical \arb{kalam} debate is not refuted but is
extended with conceptual resources unavailable to the classical
theologians.

The chapter that follows treats the unified package of Islamic
anti-concentration mechanisms (zakat, inheritance, riba
prohibition, waqf), the most macroeconomic of the worked
treatments.

\begin{summarybox}[Chapter summary]
\textbf{The doctrine}: qadar (divine decree) extends to all
events, including human actions; ikhtiyar (human choice) is real
and the basis of moral responsibility.

\textbf{The apparent contradiction}: the two doctrines together
appear to require that human choice be both determined (by qadar)
and real (for ikhtiyar). The classical theological positions
(Mu`tazili, Ash`ariyya, Maturidi) attempted resolutions but none
fully satisfying.

\textbf{The Newtonian source}: the apparent contradiction depends
on the Newtonian conception of determinism, in which determination
and motion are mutually exclusive. This conception has been
superseded by quantum mechanics.

\textbf{The Bohmian resolution}: in Bohmian mechanics, the
trajectory is fully determined by the wave function and initial
position, yet the motion is real. The structural form
``determined trajectory + real motion'' is mathematically
coherent and physically instantiated.

\textbf{The QBist resolution}: probability is perspectival.
Divine knowledge is perfect (probability 1 on the actual outcome);
human knowledge is finite (probability spread over outcomes). Both
are correct; no contradiction.

\textbf{Akbarian integration}: \arb{wahdat al-wujud} provides the
ontological substrate; \arb{tajalli} provides the dynamic content;
\arb{insan al-kamil} provides the limiting case of perfect
alignment.

\textbf{Economic implications}: effort, risk, provision under
qadar; the connection to ergodicity (agent experiences time
average, not ensemble average).

\textbf{Broader implications}: the classical kalam debate is
extended, not refuted; the modern free-will/determinism debate may
be addressing a problem whose conceptual basis has been
superseded; practical religious guidance is clarified.

\textbf{Level achieved}: Level 2. Structural problem dissolved;
operational mechanism in strong scientific sense not specified.

\textbf{What comes next}: Chapter 17 treats the unified package of
Islamic anti-concentration mechanisms.
\end{summarybox}

\cleardoublepage
